\section{\texorpdfstring{\textcolor{blue}{Point 6: plains and hills against mountains}}{Point 6: plains and hills against mountains}}\label{sec:terrain}

Does the difference between group time and solo flight depend on the terrain flown over?

\paragraph{The classes.} We class the terrain by the landforms of Hammond (\emph{Ann.\ Assoc.\ Am.\ Geogr.}\ 54, 11, 1964) in their computer version by Dikau, Brabb and Mark (\emph{USGS Open-File Report} 91-634, 1991). The elevation is the IGN BD ALTI model at \TnSource~m (Lambert-93). We average it over blocks of $\TnBlock\times\TnBlock$ cells onto a grid of \TnCell~m, a block with more than half of its cells without data is left without elevation, and so is a block whose cells all stand at exactly 0~m, which is sea delivered as elevation (\TnMaskedShare\% of the blocks with an elevation). The slope is taken by central differences on the \TnCell-m grid, and a cell is gentle when it is below \TnGentle. In a square window of $\TnWindow\times\TnWindow$ cells (\TnWindowKm~km) around a cell we take the relief, the maximum minus the minimum elevation, and the gentle share, the share of cells that are gentle. A window whose share of valid cells is below \TnCoverage{} gives the class unknown. Otherwise the cell is
\emph{plains} if the relief is below \TnPlainsRelief~m and the gentle share above \TnPlainsGentle, \emph{hills} if the relief is below \TnHillsRelief~m and it is not plains, \emph{low mountains} from \TnHillsRelief{} to \TnLowRelief~m, and \emph{high mountains} from \TnLowRelief~m up. Plains and hills make the lowland, the two mountain classes the mountains.

\paragraph{A point, a piece, an event.} A flight is its positions on the \TnGridStep-s grid of the MSD. We take each position to latitude and longitude, to Lambert-93, and read the class of the cell under it; a position whose cell has a window of low coverage, or that lies outside the model, has none. A piece takes a class when at least \TnMajority\% of its grid points are in it; otherwise it is mixed, or unknown when those points have no class. In the same flight an event takes a class when its four windows share it. Plains hold \TnPiecesPlainsG{} pieces of $\G$ and hills \TnPiecesHillsG, so in the comparison by four classes we count the plains with the hills: a piece is then hills when at least \TnMajority\% of its points are in plains or hills (\TnMergedFrom{} pieces become \TnMergedTo), and the three classes are lowland, low mountains and high mountains. Of the grid points of $\G$, \TnShareLowlandG\% lie over the lowland and \TnShareMountainG\% over the mountains, against \TnShareLowlandS\% and \TnShareMountainS\% for $\Sl$ and \TnShareLowlandA\% and \TnShareMountainA\% for $\A$.

\begin{figure}
\centering
\includegraphics[width=0.85\textwidth]{open_min0.1/terrain_forest.pdf}
\caption{$\Delta\Hloc(1000~\mathrm{s})$ by terrain class in the three designs, units in brackets: raw, $\G-\Sl$ on the pieces of dataset 1000 (lowland, low mountains, high mountains, and the two mountain classes together, open dot); same route, $(\G-\SG)-(\A-\SAA)$ on the matched pieces of $\G$; same flight, $g$ on the events. Dot: value; bar: 90\% two-way band. The band of the same-route lowland (\TnRouteLowlandDBand) extends beyond the axis.}
\label{fig:terrain-forest}
\end{figure}

\paragraph{Raw.} In the lowland $\Hk$ of $\G$ is above that of $\Sl$ by \TnRawLowlandGS{} \TnRawLowlandGSBand, in the low mountains by \TnRawLowGS{} \TnRawLowGSBand{} and in the high mountains by \TnRawHighGS{} \TnRawHighGSBand{} (Table~\ref{tab:terrain}, Figure~\ref{fig:terrain-forest}), against \RawGapHK{} \RawGapHKBand{} on all the flights (Section~\ref{sec:raw}). The gap is larger in the high mountains than in the other two classes, whose bands overlap. $\A-\Sl$ is small in the lowland (\TnRawLowlandAS{} \TnRawLowlandASBand) and in the high mountains (\TnRawHighAS{} \TnRawHighASBand), so who flies in groups explains little of the high-mountain gap: $\G-\A$ is \TnRawHighGA{} \TnRawHighGABand{} there.

\paragraph{Same flight.} On the \TnEventsLowlandClassed{} events with the four windows that lie in the lowland (\TnFlightLowlandFlights{} flights, \TnFlightLowlandPilots{} pilots) $g$ is \TnFlightLowlandG{} \TnFlightLowlandGBand, a band that includes 0; on the \TnEventsMountainClassed{} events in the mountains (\TnFlightMountainFlights{} flights, \TnFlightMountainPilots{} pilots) it is \TnFlightMountainG{} \TnFlightMountainGBand, against \PreGHK{} \PreGHKBand{} on all the \TnEventsAll{} events of Section~\ref{sec:prepost}. The two bands do not overlap. The mountains give the entry difference $D_\text{in}$ of \TnFlightMountainIn{} \TnFlightMountainInBand{} and the exit difference $D_\text{out}$ of \TnFlightMountainOut{} \TnFlightMountainOutBand; in the lowland they are \TnFlightLowlandIn{} \TnFlightLowlandInBand{} and \TnFlightLowlandOut{} \TnFlightLowlandOutBand. Events whose windows lie over several classes, \TnEventsMixedClassed{}, or over unknown ground, \TnEventsUnknownClassed, are in neither.

\paragraph{Same route.} A piece finds a match only where the pool flights fly its route: \TnRouteLowlandPiecesG{} of the \TnRawLowlandPiecesG{} pieces of $\G$ over the lowland (\TnRouteMatchedShareLowland\%) and \TnRouteMountainPiecesG{} of the \TnRawMountainPiecesG{} over the mountains (\TnRouteMatchedShareMountain\%) pass the caliper of Section~\ref{sec:matched}. We did not test why the lowland pieces rarely do. In the lowland the net gap $(\G-\SG)-(\A-\SAA)$ is \TnRouteLowlandD{} \TnRouteLowlandDBand, from \TnRouteLowlandPiecesG{} pieces: it tells nothing. In the mountains $\G-\SG$ is \TnRouteMountainG{} \TnRouteMountainGBand, $\A-\SAA$ \TnRouteMountainA{} \TnRouteMountainABand{} and the net gap \TnRouteMountainD{} \TnRouteMountainDBand{} (all the flights: \MatDeltaHKD{} \MatDeltaHKDBand), a band that includes 0.

\begin{table}
\centering\small
\begin{tabular}{@{}l r r r r l@{}}
\toprule
class & units & flights & pilots & $\Delta\Hloc(1000~\mathrm{s})$ & 90\% band \\
\midrule
\addlinespace[2pt]\multicolumn{1}{@{}l}{\emph{raw: pieces of $\G$, $\G-\Sl$}}\\
lowland & 2\,947 & 2\,355 & 593 & $+$0.010 & [$+$0.007, $+$0.014] \\
low mountains & 1\,964 & 1\,654 & 636 & $+$0.011 & [$+$0.005, $+$0.016] \\
high mountains & 8\,135 & 6\,186 & 2\,282 & $+$0.046 & [$+$0.042, $+$0.050] \\
mountains & 10\,562 & 8\,113 & 2\,640 & $+$0.027 & [$+$0.024, $+$0.030] \\
\addlinespace[2pt]\multicolumn{1}{@{}l}{\emph{same route: matched pieces of $\G$, $(\G-\SG)-(\A-\SAA)$}}\\
lowland & 13 & 13 & 12 & $-$0.003 & [$-$0.117, $+$0.111] \\
mountains & 2\,677 & 2\,435 & 1\,354 & $+$0.004 & [$-$0.007, $+$0.015] \\
\addlinespace[2pt]\multicolumn{1}{@{}l}{\emph{same flight: events, $g$}}\\
lowland & 158 & 153 & 100 & $+$0.007 & [$-$0.005, $+$0.018] \\
mountains & 1\,299 & 1\,252 & 814 & $+$0.032 & [$+$0.025, $+$0.040] \\
\bottomrule
\end{tabular}

\caption{$\Delta\Hloc(1000~\mathrm{s})$ by terrain class: pieces of $\G$ (raw, dataset 1000), matched pieces of $\G$ (same route) or events with the four windows (same flight), their flights and pilots, the difference of the design and its 90\% two-way band. Raw: $\G-\Sl$, each population's pieces taken in the class; same route: $(\G-\SG)-(\A-\SAA)$, the pieces of each side in the class; same flight: the drift-free effect $g$. The lowland of the raw design is the plains and the hills together. The pieces of the classes are in Appendix~\ref{app:terrain}.}
\label{tab:terrain}
\end{table}

\paragraph{Limits.} The class is that of the terrain under the piece or the window, not of the day or the pilot: the mountain pieces come from many more pilots (\TnRawMountainPilotsG{} against \TnRawLowlandPilotsG) who also differ in sites and in weather. The window of \TnWindowKm~km sets which relief a cell sees, and we did not try another. The comparisons condition on a class chosen at \TnMajority\%, so a piece that crosses classes is left out (\TnPiecesMixedG{} pieces of $\G$ are mixed, \TnPiecesUnknownG{} unknown). The lowland of the same flight has \TnFlightLowlandEvents{} events and that of the same route \TnRouteLowlandPiecesG{} pieces, so their bands are wide, and we read three classes and two designs at once without a correction for it.

\paragraph{What we keep.} The gain of $\Hk$ in group time rests on the mountains. In the raw design the gap of the high mountains (\TnRawHighGS{} \TnRawHighGSBand) is larger than that of the lowland (\TnRawLowlandGS{} \TnRawLowlandGSBand), and in the same flight $g$ is \TnFlightMountainG{} in the mountains against \TnFlightLowlandG{} in the lowland, with a band that includes 0. The same route says little about the lowland (\TnRouteLowlandPiecesG{} pieces) and, in the mountains, gives a net gap whose band includes 0, as it does over all the flights. The numbers show that the effect is found over the mountains; they do not say that the terrain causes it, since the pilots, the sites and the days differ between the classes.
